\documentclass[10pt,a4paper]{amsart}
\usepackage[margin=19mm]{geometry}
\usepackage{amsmath,amssymb,mathtools}
\usepackage[scr=rsfs,cal=euler]{mathalfa}
\usepackage{xfrac}
\usepackage[unicode,hidelinks,bookmarks=false]{hyperref}
\newcommand{\ie}{i.e.\ }
\newcommand{\cf}{cf.\ }
\newcommand{\resp}{respectively\ }
\newcommand{\Subsection}{\subsection}
\providecommand{\nobreakdash}{\nobreak\hbox{-}}
\setlength{\emergencystretch}{2em}
\multlinegap0pt
\usepackage{bm}
\usepackage{bbm}
\usepackage{dsfont}
\usepackage{xspace}
\usepackage{cancel}
\usepackage{mathtools}
\usepackage{amsthm}
\makeatletter
\@ifundefined{thm}{\newtheorem{thm}{Theorem}}{}
\@ifundefined{cor}{\newtheorem{cor}[thm]{Corollary}}{}
\makeatother
\makeatletter
\newcommand{\OO}{\mathcal{O}}
\expandafter\ifx\csname claim\endcsname\relax
\newenvironment{claim}[1]{\par\noindent{Claim:}\space#1}{}
\fi
\expandafter\ifx\csname claimproof\endcsname\relax
\newenvironment{claimproof}[1]{\par\noindent{Proof:}\space#1}{\hfill $\square$}
\fi
\makeatother
\title[Sums of squares of regular functions on rational surfaces]{Sums of squares of regular functions on rational surfaces}
\author{Tomasz Kowalczyk}
\date{Source: arXiv:2407.20378v2 (CC BY 4.0)}
\begin{document}
\maketitle

\noindent\emph{Source and licence.} Sums of squares of regular functions on rational surfaces. Version arXiv:2407.20378v2 (CC BY 4.0), distributed under the Creative Commons Attribution 4.0 International licence, as recorded in the arXiv licence metadata for this exact version and shown on the abstract page https://arxiv.org/abs/2407.20378v2. Full rights notice: licenses/arxiv-2407.20378-CC-BY-4.0.txt.\par\smallskip

\noindent\textit{Editorial scope.} Counted unit: the Thm whose source label is \texttt{None} in the article (source0.tex lines 460--464, source label \texttt{None}), delivered together with the article's own complete original proof of it (source0.tex lines 465--481, one \texttt{proof} environment of 17 lines). No mathematics is written, repaired or re-derived: statement and proof are the article's own text line for line. Required context retained uncounted: main (lines 335--354); cor:p(O(X))=p(O(U)) (lines 364--366). Every definition, display and label of those lines is retained unchanged. Omitted from this package and not redistributed: 1--334, 355--363, 367--459, 482--908. Nothing inside a retained excerpt is omitted or altered: every retained body is delivered exactly as the article's lines are written, so no omission receipt of any kind accompanies this package. Citations: the retained excerpts carry no citation command at all, so no bibliography entry of the article is transcribed and no reference is redistributed. Reference adaptation: none was needed and none was made; every reference inside the delivered unit resolves inside it, so no unresolved reference or citation is produced. Printed slips kept literal: no printed word, symbol, hypothesis or number of the counted statement or of its proof was altered. Layout and class: the article's own class is \texttt{amsart} with packages including geometry, amssymb, amsmath, amsfonts, inputenc, amsthm, tikz-cd, array; the wrapper is the lane's pinned \texttt{amsart} editorial wrapper at 10pt on A4, which restates the theorem-like environments the retained text needs and adds the article's own macro definitions that the retained text needs (\texttt{\textbackslash OO}), with extra wrapper packages (bm, bbm, dsfont, xspace, cancel, mathtools). The delivered numbering is the unit's own. Assets: no retained span contains a figure environment, a graphics-inclusion command, a PDF-page inclusion, a table or a TikZ picture; no image file is copied into this package, the package declares no asset dependency and no image file is redistributed with it. Licence: version arXiv:2407.20378v2 of this article is distributed under the Creative Commons Attribution 4.0 International licence, as recorded in the arXiv licence metadata for this exact version and shown on the abstract page https://arxiv.org/abs/2407.20378v2. A full rights notice is delivered at licenses/arxiv-2407.20378-CC-BY-4.0.txt. Characters: every retained excerpt is pure ASCII, so the delivered engine, XeLaTeX with its default Latin Modern text font, reports no missing character.
\begin{thm}\label{main}
Let $K$ be a formally real field. Let $C$ be an irreducible nonsingular weakly factorial curve over $K$, and $X=C\times \mathbb{A}_K$ be a product surface. Then $p(\OO(X))=p(K(X))$.

\begin{proof}
Similarly as before, since $K(X)$ is the field of fractions of $\OO(X)$, it is enough to show that the lengths of a regular functions $f$ in $K(X)$ and $\OO(X)$ are equal. By the weak factoriality of $X$, we may assume that the zero set of $f$ is finite and $f$ does not vanish on any curve. Denote the length of $f$ in $K(X)$ by $N$.


%any curve of the form $x_0 \times \mathbb{A}_K$ is given by a single equation. Hence, we can rewrite $f=f_1h^2$ where $h$ is the product of regular functions vanishing on aforementioned curves, and $f_1$ does not vanish on any such curve. We replace $f$ by $f_1$.




Of course, $K[C][y] \subset K(C)[y]$ and let $\varphi = \sum_{i=1}^N x_i^2$ be a quadratic form. By the very definition of $N$, $\varphi$ represents $f$ over $K(X)$, hence, by the Theorem of Cassels also over $K(C)[y]$. Let $f=\sum_{i=1}^N \frac{f_i^2}{g^2}$ be a presentation such that $f_i \in K[X]$ and $g \in K[C]$. If $g$ does not vanish on $C$ then we are done. Any zero $x_0 \in C$ of $g$ corresponds to a curve of the form $\{x_0 \} \times \mathbb{A}_K$ in $X$ and those are given by a single equation in $\OO(X)$ by the weak factoriality of $C$ and $X$. Hence, we can factor them from $g$ and all $f_i$ for $i=1,2,\dots, N$. However, since the factorization takes places in $\OO(X)$, it might happen that after this procedure the denominator $g$ might be replaced by a nowhere vanishing function on the whole $X$, not necessarily in $K[C]$, but in any case, $g$ does not vanish and so we obtain representation of $f$ as a sum of squares of regular functions of length $N$.

%Let $\frac{f}{g}$ be a sum of squares of rational functions of a given length $L$. Clearly, we may rewrite it as $\frac{fg}{g^2}$. We may now consider the length of $fg \in \OO(X)$ considered as a regular function or as a rational function. By the above reasoning both lengths are equal, hence we deduce the equality of Pythagoras numbers.

%By the previous lemma, the curve of the form $\{ x_0 \} \times \mathbb{A}_K$ is cut out globally by a single equation for any $x_0 \in C$. We may write $f=f_1h^2$, where the zero set of $h$ consists precisely of finite number of curves of the aforementioned type. Replacing $f$ by $f_1$ we may assume that $f$ does not not vanish on any curve of the form $\{ x_0 \} \times \mathbb{A}_K$ for any $x_0 \in C$.

\end{proof}
\end{thm}

\begin{cor}\label{cor:p(O(X))=p(O(U))}
With the assumptions as in Theorem \ref{main}. Let  $U \subset X$ be a Zariski open subset of $X$. Then $p(\OO(X))=p(\OO(U))$.
\end{cor}

\begin{thm}
Let $K$ be a formally real field. Assume that $p(K(x,y))=n < \infty$. If $X$ is a uniformly $K$-rational surface of type $m$ then the following inequality holds
$$p(\OO(X)) \leq \min\{f_2(m,n),f_4(m,n),f_8(m,n) \}. $$
In particular for large values of $m,n$ we see that $p(\OO(X))$ is at most of order $\frac{m^2n}{8}$.
\end{thm}

\begin{proof}
Fix a covering $U_i \subset X$ for $i=1,2,\dots, m$ such that each $U_i$ is isomorphic to $V_i \subset K^2$. Let $f \in \OO(X)$ be a sum of squares. Choose functions $h_i \in K[X]$ such that $\mathcal{Z}(h_i)=X\setminus U_i$ for $i=1,2,\dots, m$. Consider now $f_i:=f|_{U_i}$. By Corollary \ref{cor:p(O(X))=p(O(U))} we obtain a presentation 
$$f_i=\frac{\sum_{j=1}^n f^2_{ij}}{g_i^2}=\frac{\sum_{j=1}^n f^2_{ij}h_i^2}{g_i^2h_i^2}$$
where $\mathcal{Z}(g_i) \subset X\setminus U_i$ and $f_{ij},g_i \in K[X]$. As $\{U_i\}$ is a covering of $X$, we obtain in this manner a nowhere vanishing polynomial function $\sum_{i=1}^mg_i^2h_i^2$. By a standard trick used in real algebraic geometry, we may now add all of the numerators and denominators of all presentations to obtain a global one i.e.
$$f=\frac{\sum_{i=1}^m \sum_{j=1}^n f^2_{ij}h_i^2}{\sum_{i=1}^mg_i^2h_i^2}=\frac{\left(\sum_{i=1}^m \sum_{j=1}^n f^2_{ij}h_i^2\right)\left(\sum_{i=1}^mg_i^2h_i^2\right)}{\left(\sum_{i=1}^mg_i^2h_i^2\right)^2}.$$
This gives an upper bound on the length of $f$ as $m^2n$, in particular $p(\OO(X)) \leq m^2n$.
For the second part of the proof, recall that there are 2-, 4- and 8-square identity, in the sense of that
$(\sum_{i=1}^n x_i^2)(\sum_{i=1}^n y_i^2)=\sum_{i=1}^n z_i^2$
holds where $z_i$ is bilinear in $x_j$'s and $y_k$'s and $n=1,2,4,8$. Of course, there are such identities for every $2^k$ but the $z_i$'s are rational functions rather than bilinear in $x_j$'s and $y_k$'s.




Back to the proof, the numerator of $f$ can be written as a product of $nm$ and $m$ squares. We may now divide $nm$ and $m$ into the sets of $8$ squares each. In case $nm$ and $m$ are not divisible by $8$, we add 1 additional set, which we can fill with zeros if needed. By a repeatedly using the $8$ square identity we reduce the number of summands from $m^2n$ to $ \left\lceil\frac{nm}{8} \right\rceil \left\lceil\frac{m}{8}\right\rceil8 $ which is precisely $f_8(m,n)$. The same reasoning can be applied with the $4$-square and $2$-square identities which in total gives $p(\OO(X)) \leq \min\{f_2(m,n),f_4(m,n),f_8(m,n) \}$, finishing the proof.


\end{proof}

\end{document}
